\begin{table}[htb]
\caption{Redes neurais e referências em 30 min: tamanho, MAE de seleção e de teste, razão de generalização e custo}
\label{tab:v3-redes}
\centering
\footnotesize
\setlength{\tabcolsep}{3pt}
\renewcommand{\arraystretch}{1.15}
\begin{tabular}{|l|c|c|c|c|c|c|}
\hline
\textbf{Modelo} & \textbf{Tamanho} & \textbf{MAE seleção} & \textbf{MAE teste} & \textbf{Razão norm.} & \textbf{Reajuste (s)} & \textbf{CUDA (MiB)} \\
\hline
DLinear & 78 & 3{,}91 & 5{,}14 & 1{,}16 & 0{,}813 & 65 \\
TCN & 9793 & 3{,}81 & 5{,}38 & 1{,}24 & 0{,}501 & 68 \\
N-BEATSx & 40576 & 3{,}93 & 5{,}25 & 1{,}18 & 0{,}269 & 66 \\
N-HiTS & 31983 & 3{,}91 & 5{,}73 & 1{,}29 & 0{,}305 & 66 \\
PatchTST & 18145 & 3{,}99 & 4{,}32 & 0{,}95 & 0{,}373 & 73 \\
XGBoost (V2) & 2789 folhas & 3{,}80 & 4{,}33 & 0{,}98 & 0{,}075 & -- \\
Kalman + XGBoost & 793 folhas & 3{,}88 & 4{,}32 & 0{,}98 & 0{,}118 & -- \\
\hline
\end{tabular}\\
\footnotesize Razão norm. = (teste/seleção) dividido pela mesma razão da persistência. Redes: CUDA, corrida original; XGBoost e híbrido: CPU, medido na análise.\\
\textbf{\footnotesize Fonte: Autoria própria.}
\end{table}
